\documentclass[a4paper,12pt]{article}
\usepackage[utf8]{inputenc}
\usepackage{babel}
\usepackage[T1]{fontenc}
\usepackage{enumitem}
\usepackage{geometry}
\usepackage{listings}
\usepackage{xcolor}
\usepackage{tcolorbox}
\usepackage{hyperref}
\usepackage{multicol}
\usepackage{array}
\usepackage{booktabs}
\usepackage{tabularx}
\usepackage{amssymb}

\geometry{margin=2.5cm}
\title{Tutorato Programmazione 2}
\author{Michele Porcellato, Giacomo Vettore}
\date{13 Maggio 2026}

% Java listing style
\lstdefinestyle{java}{
  language=Java,
  basicstyle=\ttfamily\small,
  keywordstyle=\color{blue}\bfseries,
  commentstyle=\color{gray},
  stringstyle=\color{orange!80!black},
  numberstyle=\ttfamily\scriptsize\color{gray},
  numbers=left,
  numbersep=8pt,
  frame=single,
  frameround=tttt,
  rulecolor=\color{black!30},
  backgroundcolor=\color{gray!5},
  breaklines=true,
  tabsize=2,
  showstringspaces=false,
  xleftmargin=1.5em,
}

% Box for error answers
\newcommand{\errorbox}{%
  \vspace{0.4em}
  \noindent
  \begin{tabular}{|p{3.2cm}|p{3.4cm}|p{5.6cm}|}
    \hline
    \rule{0pt}{2.4em}errore alla riga n.\dotfill\dotfill &
    \textbf{A} in compilazione\newline\textbf{B} a runtime &
    perché $\rightarrow$ \\[1.6em]
    \hline
  \end{tabular}
  \vspace{0.4em}
}

% Box for output answers
\newcommand{\outputbox}[1]{%
  \vspace{0.4em}
  \noindent
  \begin{tabular}{|l|p{10cm}|}
    \hline
    \rule{0pt}{1.8em}\textbf{#1} \quad l'output è $\rightarrow$ & \phantom{XXXXXXXXXXXXXXXXXXX} \\
    \hline
  \end{tabular}
  \vspace{0.4em}
}

\begin{document}

\maketitle
\thispagestyle{empty}

% ─── Istruzioni ───────────────────────────────────────────────────────────────
\begin{tcolorbox}[colback=gray!8, colframe=black!60, title=\textbf{Istruzioni (leggere attentamente)}]
Vi sono tre tipi di domande e, di conseguenza, risposte:

\begin{enumerate}[leftmargin=*, label=\arabic*.]

\item Un frammento di codice che esegue correttamente. Si indichi l'output nell'apposito riquadro.

\smallskip
\noindent
\begin{tabular}{|l|p{8cm}|}
\hline
\rule{0pt}{1.6em}\textbf{Test N} \quad l'output è $\rightarrow$ & \\
\hline
\end{tabular}

\item Un frammento di codice che genera errori. Si indichi in quale riga l'errore avviene; si specifichi la tipologia (in compilazione o a runtime) cerchiando la cella corrispondente (\textbf{A} o \textbf{B}); si spieghi brevemente la ragione dell'errore. \emph{Tutti e tre i riquadri devono essere compilati.}

\smallskip
\noindent
\begin{tabularx}{\linewidth}{|p{3.2cm}|p{4cm}|X|}
\hline
\rule{0pt}{2.2em}errore alla riga n.\dotfill &
\begin{tabular}{@{}l@{}}
  \textbf{A} in compilazione \\[2pt]
  \textbf{B} a runtime
\end{tabular} &
perché $\rightarrow$ \\
\hline
\end{tabularx}

\item Domande vero/falso. Si riporti \textbf{V} (vero) o \textbf{F} (falso) nelle caselle corrispondenti. In questa tipologia, le \emph{risposte errate sottraggono punti}.

\smallskip
\noindent
\begin{tabular}{|c|c|c|c|c|c|c|c|}
\hline
\rule{0pt}{1.4em}\textbf{A} & \textbf{B} & \textbf{C} & \textbf{D} & \textbf{E} & \textbf{F} & \textbf{G} & \textbf{H} \\[2pt]
\hline
\phantom{XX} & \phantom{XX} & \phantom{XX} & \phantom{XX} & \phantom{XX} & \phantom{XX} & \phantom{XX} & \phantom{XX} \\
\hline
\end{tabular}

\end{enumerate}
\end{tcolorbox}



\newpage

% ═══════════════════════════════════════════════════════════════════════════════
%  TEST 1  — ERRORE
% ═══════════════════════════════════════════════════════════════════════════════
\section*{Test 1}
\errorbox

\begin{lstlisting}[style=java]
public class Test {
    static final int MAX = 10;
    public static void main(String[] args) {
        A a = new A();
        a.m1(1);
        C c = new C();
        c.m2(2);
    }
}
interface I1 {
    public int m1(int x);
}
interface I2 extends I1 {
    public int m2(int x);
}
class A implements I1 {
    int a = 20;
    public int m1(int x) {
        return a + x;
    }
}
class B extends A {}
class C extends B implements I2 {}
\end{lstlisting}

% ═══════════════════════════════════════════════════════════════════════════════
%  TEST 2  — ERRORE
% ═══════════════════════════════════════════════════════════════════════════════
\section*{Test 2}
\errorbox

\begin{lstlisting}[style=java]
public class Test {
    static final int MAX = 10;
    public static void main(String[] args) {
        A a = new A(10);
        a.m("mela");
        A a1 = new C();
        C c1 = (C) a1;
        c1.m("pera");
        A a2 = new B();
        C c2 = (C) a2;
        c2.m("banana");
    }
}
interface I1 {
    public int m(String s);
}
class A implements I1 {
    int x;
    A(int x) { this.x = x; }
    public int m(String s) {
        return s.length() * x;
    }
}
class B extends A {
    B() { super(0); }
}
class C extends B {}
\end{lstlisting}

% ═══════════════════════════════════════════════════════════════════════════════
%  TEST 3  — ERRORE
% ═══════════════════════════════════════════════════════════════════════════════
\section*{Test 3}
\errorbox

\begin{lstlisting}[style=java]
public class Test {
    public static void main(String[] args) {
        A obj = new B();
        obj.m(new D());
    }
}
class A {
    final void m(C c) { System.out.println("1"); }
}
class B extends A {
    void m(C c) { System.out.println("2"); }
    void m(D c) { System.out.println("3"); }
}
class C {}
class D extends C {}
\end{lstlisting}

% ═══════════════════════════════════════════════════════════════════════════════
%  TEST 4  — OUTPUT
% ═══════════════════════════════════════════════════════════════════════════════
\section*{Test 4}
\outputbox{Test 4}

\begin{lstlisting}[style=java]
public class Test {
    public static void main(String[] args) {
        A obj = new B();
        A par = new B();
        System.out.println(obj.m(par));
    }
}
class A {
    int x = 0;
    String m(A a) { return "a in a"; }
}
class B extends A {
    String m(B b) { return "b in b"; }
    String m(A a) { return "a in b"; }
}
\end{lstlisting}

% ═══════════════════════════════════════════════════════════════════════════════
%  TEST 5  — OUTPUT
% ═══════════════════════════════════════════════════════════════════════════════
\section*{Test 5}
\outputbox{Test 5}

\begin{lstlisting}[style=java]
public class Test {
    public static void main(String[] args) {
        A a1 = new A(5);
        A a2 = new A(3);
        A a3 = new A(5);
        System.out.println(a1.equals(a2));
        System.out.println(!a1.equals(a3));
    }
}
class A {
    int x = 0;
    A(int x) { this.x = x; }
}
\end{lstlisting}

% ═══════════════════════════════════════════════════════════════════════════════
%  TEST 6  — OUTPUT
% ═══════════════════════════════════════════════════════════════════════════════
\section*{Test 6}
\outputbox{Test 6}

\begin{lstlisting}[style=java]
public class Test {
    public static void main(String[] args) {
        A a = new A();
        int r1 = a.m(6);
        A a2 = new B(new Z());
        int r2 = a2.m(1);
        B b = new B(new Z());
        int r3 = b.m(2);
        System.out.println("" + r1 + " " + r2 + " " + r3);
    }
}
class Z {}
class A {
    static int val = 5;
    int m(int x) { val--; return x * val; }
}
class B extends A {
    Z z;
    B(Z z) { this.z = z; val++; }
    int m(int x) { return x * val + 1; }
}
\end{lstlisting}

% ═══════════════════════════════════════════════════════════════════════════════
%  TEST 7  — OUTPUT
% ═══════════════════════════════════════════════════════════════════════════════
\section*{Test 7}
\outputbox{Test 7}

\begin{lstlisting}[style=java]
public class Test {
    public static void main(String[] args) {
        I i = new C(4);
        System.out.println(i.m(2));
    }
}
interface I {
    int m(int z);
}
class A implements I {
    int x;
    A(int x) { this.x = x + 1; }
    public int m(int z) { return x + z; }
}
class B extends A {
    B(int x) { super(++x); }
    public int m(int z) { return x * z; }
}
class C extends B {
    C(int x) { super(++x); }
}
\end{lstlisting}

% ═══════════════════════════════════════════════════════════════════════════════
%  TEST 8  — OUTPUT
% ═══════════════════════════════════════════════════════════════════════════════
\section*{Test 8}
\outputbox{Test 8}

\begin{lstlisting}[style=java]
import java.util.*;
public class Test {
    public static void main(String[] args) {
        String[] a = {"limone", "ananas", "mango", "lime"};
        Set<Integer> s = new HashSet<>();
        List<Integer> l = new ArrayList<>();
        for (String i : a) {
            s.add(i.length());
            l.add(i.length());
        }
        System.out.println(s.size() + l.size());
    }
}
\end{lstlisting}

% ═══════════════════════════════════════════════════════════════════════════════
%  TEST 9  — VERO / FALSO
% ═══════════════════════════════════════════════════════════════════════════════
\section*{Test 9}

\noindent
\begin{tabular}{|c|c|c|c|c|c|c|c|}
\hline
\rule{0pt}{1.4em}\textbf{A} & \textbf{B} & \textbf{C} & \textbf{D} & \textbf{E} & \textbf{F} & \textbf{G} & \textbf{H} \\[2pt]
\hline
\phantom{XXX} & \phantom{XXX} & \phantom{XXX} & \phantom{XXX} & \phantom{XXX} & \phantom{XXX} & \phantom{XXX} & \phantom{XXX} \\[8pt]
\hline
\end{tabular}

\bigskip

\begin{enumerate}[label=\textbf{\Alph*.}, leftmargin=*, itemsep=0.5em]

\item In una classe possono coesistere più metodi con lo stesso nome ma firme
      (numero e tipo dei parametri) diverse.

\item La parola chiave \texttt{finally} serve a terminare l'esecuzione del programma.

\item Si consideri un attributo \texttt{x} dichiarato come \texttt{protected} nella classe
      \texttt{C} del package \texttt{P1}; \texttt{x} non è visibile da una classe \texttt{D}
      appartenente a un package \texttt{P2}, a meno che \texttt{D} non erediti da \texttt{C}.

\item Le istruzioni \texttt{Object[] x = new Object[10]; x[0] = "elemento";} sono
      illegali in Java, perché per gli array non è supportato il polimorfismo; questo è il
      motivo principale per l'esistenza del Java Collection Framework, che invece supporta
      questa funzionalità.

\item Il garbage collector è una funzionalità del supporto run-time di Java molto utile,
      in quanto gestisce in maniera automatica la deallocazione di oggetti non più utilizzati
      dal programma, semplificando il lavoro del programmatore ed evitandone potenziali errori.

\item Quando in Java si parla di un metodo \emph{overloaded} si fa riferimento a un metodo
      ``sovraccarico'', cioè la cui implementazione genera una computazione particolarmente pesante.

\item Una classe Java definita come \texttt{abstract} può essere usata all'interno di
      gerarchie di classi con ereditarietà multipla.

\item Siano dati due oggetti \texttt{a} e \texttt{b}, per i quali il test
      \texttt{a.hashCode() == b.hashCode()} ritorna \texttt{true}. In questo caso, si può
      stabilire con certezza che \texttt{a} e \texttt{b} sono identici, cioè puntano allo
      stesso oggetto.

\end{enumerate}

% ═══════════════════════════════════════════════════════════════════════════════
%  TEST 10  — OUTPUT
% ═══════════════════════════════════════════════════════════════════════════════
\section*{Test 10}
\outputbox{Test 10}
 
\begin{lstlisting}[style=java]
public class Test {
    public static void main(String[] args){
        A obj = new B();
        obj.m(new D());
    }
}
class A {
    void m(C c) { System.out.println("x"); }
}
class B extends A {
    void m(C c) { System.out.println("y"); }
    void m(D d) { System.out.println("z"); }
}
class C { }
class D extends C { }
\end{lstlisting}
 
% ═══════════════════════════════════════════════════════════════════════════════
%  TEST 11  — ERRORE
% ═══════════════════════════════════════════════════════════════════════════════
\section*{Test 11}
\errorbox
 
\begin{lstlisting}[style=java]
import java.util.*;
public class Sette {
    Sette(){
        List a = new ArrayList();
        List b = new HashSet();
        for (int k=0; k<10; k++) {
            String s = "A" + (k%3);
            a.add(s);
            b.add(s);
        }
        System.out.println(a.size() + " " + b.size());
    }
    public static void main(String[] a) { new Sette(); }
    public static void main(String a)   { new Sette(); new Sette(); }
}
\end{lstlisting}
 
% ═══════════════════════════════════════════════════════════════════════════════
%  TEST 12  — OUTPUT
% ═══════════════════════════════════════════════════════════════════════════════
\section*{Test 12}
\outputbox{Test 12}
 
\begin{lstlisting}[style=java]
public class Test {
    public static void main(String[] args) {
        A a = new B();
        B b = new B();
        J j = b;
        if (a.equals(b))
            System.out.println(j.m(1) + b.m("scritto"));
        else System.out.println(j.m(3));
    }
}
interface I { int m(int z); }
interface J extends I {}
class A implements I {
    int x = 10;
    public int m(int z) { return x + z; }
    public boolean equals(Object o) {
        A obj = (A) o;
        return x == obj.x;
    }
}
class B extends A implements J {
    public int m(String s) { return s.length(); }
}
\end{lstlisting}
 
% ═══════════════════════════════════════════════════════════════════════════════
%  TEST 13  — OUTPUT
% ═══════════════════════════════════════════════════════════════════════════════
\section*{Test 13}
\outputbox{Test 13}
 
\begin{lstlisting}[style=java]
import java.util.*;
public class E6 {
    static Collection ll = new LinkedList();
    int x = 6;
    E6() {}
    E6(int x) {
        ll.add(this);
        ll.add(new EA5());
    }
    public static void main(String z[]) {
        new E6(3);
        Iterator iter = ll.iterator();
        while (iter.hasNext()) {
            ((E6)(iter.next())).f();
        }
    }
    public void f() { System.out.print(x); }
}
class EA5 extends E6 {
    public void f() {
        x++; super.f(); System.out.print(2);
    }
}
\end{lstlisting}
 
% ═══════════════════════════════════════════════════════════════════════════════
%  TEST 14  — OUTPUT
% ═══════════════════════════════════════════════════════════════════════════════
\section*{Test 14}
\outputbox{Test 14}
 
\begin{lstlisting}[style=java]
public class Sei {
    char f() { return '6'; }
    public static void main(String e[]) {
        Sei a = new Sei();
        Sei b = new Sette();
        Sette c = new Sette();
        System.out.print(a.f() + " " + b.f() + " " + c.f() + " ");
        char ch[] = {'A', 'B', 'A', 'B', 'A', 'B'};
        int i1 = 0, i2 = 2, i3 = 4;
        if (a.equals(b)) i1++;
        if (b.equals(a)) i2++;
        if (c.equals(b)) i3++;
        System.out.print(ch[i1] + " " + ch[i2] + " " + ch[i3]);
    }
}
class Sette extends Sei {
    char f() { return '7'; }
    public boolean equals(Object a) { return (a instanceof Sei); }
    public int hashCode() { return 3; }
}
\end{lstlisting}
 
% ═══════════════════════════════════════════════════════════════════════════════
%  TEST 15  — ERRORE
% ═══════════════════════════════════════════════════════════════════════════════
\section*{Test 15}
\errorbox
 
\begin{lstlisting}[style=java]
public class Tre {
    class A {
        public A(int k) { System.out.print(k); }
        public void finalize() { System.out.print("A"); }
    }
    class B extends A {
        public B(int k) { System.out.print(k); }
        public void finalize() { System.out.print("A"); }
    }
    public static void main(String z[]) { new Tre(); }
    Tre() {
        A a = new B(3);
        B b = (B) a;
        a = null;
        b = new B(3);
        System.gc();
        System.runFinalization();
    }
}
\end{lstlisting}
 
% ═══════════════════════════════════════════════════════════════════════════════
%  TEST 16  — OUTPUT
% ═══════════════════════════════════════════════════════════════════════════════
\section*{Test 16}
\outputbox{Test 16}
 
\begin{lstlisting}[style=java]
public class Test8 {
    public static void main(String[] args) {
        A al = new A(5);
        A a2 = new A(al.m()/2);
        boolean b = (al == a2);
        System.out.print(b);
    }
}
class A {
    int x;
    A(int x) { this.x = x; }
    int m() { return x*2; }
}
\end{lstlisting}
 
% ═══════════════════════════════════════════════════════════════════════════════
%  TEST 17  — ERRORE
% ═══════════════════════════════════════════════════════════════════════════════
\section*{Test 17}
\errorbox
 
\begin{lstlisting}[style=java]
import java.util.*;
public class B {
    B() {
        Collection b = new Collection();
        for (int k=0; k<10; k++) {
            String s = "A" + (k%4);
            b.add(s);
        }
        int count = 0;
        Iterator i = b.iterator();
        while (i.hasNext()) {
            Object s = i.next();
            count++;
        }
        System.out.println(count);
    }
    public static void main(String[] a) { new B(); new B(); }
    public static void main(String a)   { new B(); }
}
\end{lstlisting}
 
% ═══════════════════════════════════════════════════════════════════════════════
%  TEST 18  — OUTPUT
% ═══════════════════════════════════════════════════════════════════════════════
\section*{Test 18}
\outputbox{Test 18}
 
\begin{lstlisting}[style=java]
public class C {
    public static int x;
    C(int s) { x = s; }
    void f() { System.out.print(x); }
    public static void main(String a[]) {
        C b = new C(3);
        C c = new C(5);
        b.f();
        c.f();
    }
}
\end{lstlisting}
 
\end{document}